\documentclass[aspectratio=169, 10pt]{beamer}

% --- Pacotes de Idioma e Codificacao ---
\usepackage[utf8]{inputenc}
\usepackage[T1]{fontenc}
\usepackage[brazil]{babel}
\usepackage{booktabs}
\usepackage{graphicx}
\usepackage{tikz}
\usepackage{hyperref}
\usetikzlibrary{shapes.geometric, arrows.meta, positioning}

% --- Configuracao de Tema e Cores ---
\usetheme{Madrid}
\usecolortheme{whale}

% Paleta de cores corporativa e sobria
\definecolor{primaryNavy}{RGB}{20, 35, 60}
\definecolor{accentTeal}{RGB}{0, 130, 120}
\definecolor{lightGray}{RGB}{245, 247, 250}

\setbeamercolor{palette primary}{bg=primaryNavy, fg=white}
\setbeamercolor{palette secondary}{bg=accentTeal, fg=white}
\setbeamercolor{palette tertiary}{bg=primaryNavy!90, fg=white}
\setbeamercolor{structure}{fg=primaryNavy}
\setbeamercolor{titlelike}{parent=palette primary}
\setbeamercolor{block title}{bg=primaryNavy, fg=white}
\setbeamercolor{block body}{bg=lightGray, fg=black}
\setbeamercolor{block title alerted}{bg=red!75!black, fg=white}
\setbeamercolor{block body alerted}{bg=red!5, fg=black}

\setbeamertemplate{navigation symbols}{}

% --- Metadados da Apresentacao ---
\title[HelpDesk AI — NexusTech]{\textbf{HelpDesk AI: Triagem Inteligente por Setor Corporativo}}
\subtitle{Estudo de Caso NexusTech S.A. — No Code, Low Code e Vibecode}

\author[Gustavo Negreiros]{
    \textbf{Discente:} Gustavo Negreiros
}

\institute[UDF]{
    \textbf{Centro Universitário do Distrito Federal (UDF)} \\
    Tendências em Ciência da Computação \\
    \textbf{Docente:} Profa. Kadidja Valéria
}

\date{\today}

% --- Inicio do Documento ---
\begin{document}

% -----------------------------------------------------------------------------
% Slide 1: Capa
% -----------------------------------------------------------------------------
\begin{frame}[plain]
    \titlepage
\end{frame}

% -----------------------------------------------------------------------------
% Slide 2: Cenario Corporativo e O Desafio
% -----------------------------------------------------------------------------
\begin{frame}{Cenário Corporativo: NexusTech S.A.}
    \begin{columns}[T]
        \begin{column}{0.48\textwidth}
            \begin{block}{O Contexto Organizacional}
                \begin{itemize}
                    \item Empresa de médio porte de tecnologia financeira e logística com \textbf{5 departamentos}.
                    \item Fila única de chamados sem distinção de criticidade entre setores.
                    \item Analistas sobrecarregados com triagem manual e dúvidas recorrentes de Nível 1.
                \end{itemize}
            \end{block}
        \end{column}
        
        \begin{column}{0.48\textwidth}
            \begin{alertblock}{O Problema da Fila Única}
                \begin{itemize}
                    \item Uma falha simples no setor Administrativo competia na mesma fila que um incidente crítico em \textbf{Operações (NOC)} ou \textbf{Faturamento}.
                    \item Subjetividade: usuários marcavam todos os chamados como urgentes.
                \end{itemize}
            \end{alertblock}
        \end{column}
    \end{columns}
\end{frame}

% -----------------------------------------------------------------------------
% Slide 3: Matriz de Criticidade por Setor
% -----------------------------------------------------------------------------
\begin{frame}{Estrutura Organizacional e Matriz de Criticidade}
    \begin{table}[h]
        \centering
        \footnotesize
        \begin{tabular}{llll}
            \toprule
            \textbf{Setor (NexusTech S.A.)} & \textbf{Impacto de Negócio} & \textbf{Exemplos de Falhas} & \textbf{Prioridade / SLA} \\
            \midrule
            \textbf{1. Operações e NOC} & Altíssimo (Cliente Final) & Queda de servidores / Links & P1 ou P2 (15m -- 1h) \\
            \textbf{2. Financeiro e Faturamento} & Alto (Fluxo de Caixa) & Erro em ERP / Notas Fiscais & P1 ou P2 (30m -- 2h) \\
            \textbf{3. Engenharia de Software} & Médio/Alto (Deploys) & Falha em CI/CD / Acesso Git & P2 ou P3 (1h -- 4h) \\
            \textbf{4. Comercial e CX} & Médio (Vendas e CRM) & Instabilidade em CRM / VoIP & P2 ou P3 (2h -- 6h) \\
            \textbf{5. RH e Administrativo} & Moderado (Interno) & Impressoras / Contas Office & P3 ou P4 (4h -- 24h) \\
            \bottomrule
        \end{tabular}
    \end{table}

    \vspace{0.2cm}
    \begin{block}{Regra de Negócio Contextual}
        A prioridade final e o SLA não dependem apenas do sintoma técnico, mas do \textbf{cruzamento entre a gravidade do relato e o peso estratégico do setor solicitante}.
    \end{block}
\end{frame}

% -----------------------------------------------------------------------------
% Slide 4: A Solucao: HelpDesk AI Contextual
% -----------------------------------------------------------------------------
\begin{frame}{A Solução: HelpDesk AI}
    \begin{block}{\textbf{Objetivo do Protótipo}}
        Automatizar a triagem primária e o autoatendimento da NexusTech S.A. cruzando o \textbf{Setor Corporativo} e o \textbf{Relato Técnico} por meio de plataformas visuais e IA generativa.
    \end{block}

    \vspace{0.3cm}
    \textbf{Entregas Automatizadas pelo Sistema:}
    \begin{itemize}
        \item \textbf{Classificação Técnica:} Hardware, Software, Rede/Infraestrutura ou Acesso/Segurança.
        \item \textbf{Cálculo Dinâmico de Prioridade e SLA:} Atribuição racional de P1 a P4 e tempo estimado de atendimento conforme o setor.
        \item \textbf{Definição de Fila:} Direcionamento para Nível 1, Nível 2 Infra ou Nível 2 Sistemas.
        \item \textbf{Autoatendimento Imediato:} Envio de guia passo a passo por e-mail para resolução rápida pelo próprio colaborador.
    \end{itemize}
\end{frame}

% -----------------------------------------------------------------------------
% Slide 5: O Tripe Metodologico
% -----------------------------------------------------------------------------
\begin{frame}{O Tripé Metodológico Aplicado}
    \begin{columns}[T]
        \begin{column}{0.32\textwidth}
            \begin{block}{\textbf{1. No Code}}
                \textbf{Google Forms \& Sheets}
                \begin{itemize}
                    \item Formulário com seleção de Setor Corporativo.
                    \item Planilha central de governança de TI.
                    \item Zero linhas de código.
                \end{itemize}
            \end{block}
        \end{column}

        \begin{column}{0.32\textwidth}
            \begin{block}{\textbf{2. Low Code}}
                \textbf{Make (Integromat)}
                \begin{itemize}
                    \item Escuta ativa de novos chamados.
                    \item Injeção das variáveis de Setor e Relato na API.
                    \item Atualização da planilha e disparo de e-mail.
                \end{itemize}
            \end{block}
        \end{column}

        \begin{column}{0.32\textwidth}
            \begin{block}{\textbf{3. Vibecode}}
                \textbf{Gemini / LLM}
                \begin{itemize}
                    \item Avaliação contextual (Setor $\times$ Incidente).
                    \item Definição de SLA e Fila Técnica.
                    \item Redação de guia de autoatendimento.
                \end{itemize}
            \end{block}
        \end{column}
    \end{columns}
\end{frame}

% -----------------------------------------------------------------------------
% Slide 6: Camadas No Code e Low Code na Pratica
% -----------------------------------------------------------------------------
\begin{frame}{Implementação: No Code e Low Code}
    \begin{columns}[T]
        \begin{column}{0.48\textwidth}
            \textbf{No Code (Portal e Governança):}
            \begin{itemize}
                \item \textbf{Entrada:} Nome, E-mail, \textbf{Setor (Lista Suspensa)}, Título e Descrição.
                \item \textbf{Colunas de Governança no Sheets:} Categoria, Prioridade, SLA Estimado e Fila Responsável (preenchidas automaticamente).
            \end{itemize}
        \end{column}

        \begin{column}{0.48\textwidth}
            \textbf{Low Code (Orquestração no Make):}
            \begin{itemize}
                \item Captura a linha recém-criada no Google Sheets.
                \item Envia os dados estruturados para o modelo Gemini.
                \item Grava o diagnóstico na planilha da TI e notifica o colaborador via Gmail.
            \end{itemize}
        \end{column}
    \end{columns}

    \vspace{0.4cm}
    \begin{center}
        \begin{tikzpicture}[node distance=2.2cm, auto, >=Stealth]
            \tikzstyle{box} = [rectangle, draw=primaryNavy, fill=white, thick, text centered, rounded corners, minimum height=0.9cm, minimum width=2.5cm]
            \node[box] (n1) {\textbf{Ticket + Setor}};
            \node[box, right of=n1, xshift=1.2cm] (n2) {\textbf{Make Trigger}};
            \node[box, right of=n2, xshift=1.2cm] (n3) {\textbf{Matriz IA}};
            \node[box, right of=n3, xshift=1.2cm] (n4) {\textbf{Sheets + Gmail}};
            
            \draw[->, thick, color=accentTeal] (n1) -- (n2);
            \draw[->, thick, color=accentTeal] (n2) -- (n3);
            \draw[->, thick, color=accentTeal] (n3) -- (n4);
        \end{tikzpicture}
    \end{center}
\end{frame}

% -----------------------------------------------------------------------------
% Slide 7: Camada Vibecode: Matriz Cognitiva
% -----------------------------------------------------------------------------
\begin{frame}{Camada Vibecode: Ponderação Setor $\times$ Incidente}
    \textbf{Por que usar Vibecode em vez de condicionais tradicionais?}
    \begin{itemize}
        \item Em programação tradicional, cruzar 5 setores com centenas de sintomas diferentes exigiria árvores complexas de \texttt{if/else} ou expressões regulares frágeis.
        \item Com engenharia de prompt, a IA compreende a semântica do problema sob o contexto do setor.
    \end{itemize}

    \vspace{0.2cm}
    \begin{exampleblock}{Estrutura do Prompt Contextual (Vibecode)}
        \textit{"Você é o motor de triagem de TI da NexusTech S.A. Analise o chamado aberto pelo setor \{\{Setor\}\} com o título \{\{Título\}\} e relato: \{\{Descrição\}\}. Considerando o peso estratégico desse setor, retorne: 1) Categoria Técnica, 2) Prioridade (P1 a P4) e SLA estimado, 3) Fila de Escalonamento e 4) Um guia prático em 3 passos para autoatendimento imediato do colaborador."}
    \end{exampleblock}
\end{frame}

% -----------------------------------------------------------------------------
% Slide 8: Arquitetura Ponta a Ponta
% -----------------------------------------------------------------------------
\begin{frame}{Arquitetura e Fluxo de Dados Ponta a Ponta}
    \begin{center}
        \begin{tikzpicture}[node distance=1.6cm, auto, >=Stealth, font=\small]
            \tikzstyle{step} = [rectangle, draw=primaryNavy, fill=primaryNavy!10, thick, text centered, rounded corners, minimum height=1.2cm, text width=2.2cm]
            
            \node[step] (f1) {\textbf{1. Colaborador}\\Informa Setor e Falha};
            \node[step, right of=f1, xshift=1.1cm] (f2) {\textbf{2. Sheets}\\Registra Ticket};
            \node[step, right of=f2, xshift=1.1cm] (f3) {\textbf{3. Make}\\Extrai Contexto};
            \node[step, right of=f3, xshift=1.1cm] (f4) {\textbf{4. Gemini IA}\\Calcula SLA e Guia};
            \node[step, right of=f4, xshift=1.1cm] (f5) {\textbf{5. Desfecho}\\Atualiza TI + E-mail};

            \draw[->, very thick, color=accentTeal] (f1) -- (f2);
            \draw[->, very thick, color=accentTeal] (f2) -- (f3);
            \draw[->, very thick, color=accentTeal] (f3) -- (f4);
            \draw[->, very thick, color=accentTeal] (f4) -- (f5);
        \end{tikzpicture}
    \end{center}
    
    \vspace{0.4cm}
    \begin{block}{Exemplo Prático de Comportamento do Sistema}
        \begin{itemize}
            \item \textbf{Caso A:} "Lentidão na planilha" no setor \textbf{RH} $\rightarrow$ Classificado como \textbf{P4 (SLA 24h)}.
            \item \textbf{Caso B:} "Erro ao emitir nota fiscal no ERP" no setor \textbf{Financeiro} $\rightarrow$ Classificado como \textbf{P1 (SLA 30 min)} e escalonado para Nível 2 Sistemas.
        \end{itemize}
    \end{block}
\end{frame}

% -----------------------------------------------------------------------------
% Slide 9: Reflexao Critica: 3 Vantagens
% -----------------------------------------------------------------------------
\begin{frame}{Reflexão Crítica: 3 Vantagens Principais}
    \begin{enumerate}
        \item \textbf{Priorização Inteligente Sensível ao Negócio (\textit{Business-Aware Triage}):}
        \begin{itemize}
            \item Elimina a subjetividade do usuário ao ponderar automaticamente a criticidade do departamento afetado na NexusTech S.A.
        \end{itemize}
        \vspace{0.2cm}
        \item \textbf{Redução Drástica do Tempo de Resposta em Nível 1:}
        \begin{itemize}
            \item O colaborador recebe instruções técnicas personalizadas e a previsão de SLA em segundos após submeter o formulário.
        \end{itemize}
        \vspace{0.2cm}
        \item \textbf{Alta Velocidade de Prototipação e Custo Zero:}
        \begin{itemize}
            \item Arquitetura completa validada sem necessidade de provisionamento de servidores, bancos relacionais ou código boilerplate.
        \end{itemize}
    \end{enumerate}
\end{frame}

% -----------------------------------------------------------------------------
% Slide 10: Reflexao Critica: 3 Limitacoes
% -----------------------------------------------------------------------------
\begin{frame}{Reflexão Crítica: 3 Limitações e Desafios}
    \begin{enumerate}
        \item \textbf{Risco de Diagnóstico Incompleto ou Alucinação:}
        \begin{itemize}
            \item Relatos muito curtos ou ambíguos podem induzir a IA a sugerir passos genéricos. O sistema atua como triagem inicial, mantendo a supervisão humana na planilha.
        \end{itemize}
        \vspace{0.2cm}
        \item \textbf{Limites de Operações e Latência no Plano Gratuito:}
        \begin{itemize}
            \item A conta gratuita do Make possui limite mensal de 1.000 operações e janelas de verificação periódica (polling).
        \end{itemize}
        \vspace{0.2cm}
        \item \textbf{Governança, Segurança e Privacidade (LGPD):}
        \begin{itemize}
            \item Dados corporativos da NexusTech S.A. transitam por APIs públicas. Em produção real, exige-se anonimização e instâncias empresariais privadas.
        \end{itemize}
    \end{enumerate}
\end{frame}

% -----------------------------------------------------------------------------
% Slide 11: Proximos Passos e Conclusao
% -----------------------------------------------------------------------------
\begin{frame}{Próximos Passos e Conclusão}
    \begin{columns}[T]
        \begin{column}{0.48\textwidth}
            \textbf{Evoluções Futuras na NexusTech S.A.:}
            \begin{itemize}
                \item Roteamento de alertas críticos (P1) diretamente para canais do Slack ou Microsoft Teams da equipe de plantão.
                \item Uso de RAG (\textit{Retrieval-Augmented Generation}) conectado aos manuais internos da empresa.
                \item Integração bidirecional com Jira Service Management.
            \end{itemize}
        \end{column}

        \begin{column}{0.48\textwidth}
            \begin{block}{\textbf{Conclusão}}
                A combinação de No Code, Low Code e Vibecode permitiu criar uma solução de triagem corporativa contextualizada, demonstrando como a IA generativa eleva a eficiência operacional na Engenharia de Software moderna.
            \end{block}
        \end{column}
    \end{columns}

    \vspace{0.5cm}
    \begin{center}
        \large \textbf{Obrigado!} \\
        \normalsize Espaço aberto para dúvidas e considerações.
    \end{center}
\end{frame}

\end{document}
